\documentclass[a4paper,10pt]{article}

\usepackage{latexsym}
\usepackage[empty]{fullpage}
\usepackage{titlesec}
\usepackage{marvosym}
\usepackage[usenames,dvipsnames]{color}
\usepackage{verbatim}
\usepackage{enumitem}
\usepackage[pdftex]{hyperref}
\usepackage{fancyhdr}
\usepackage{mathptmx}
\usepackage{csquotes}

\pagestyle{fancy}
\fancyhf{}
\renewcommand{\headrulewidth}{0pt}
\renewcommand{\footrulewidth}{0pt}

\addtolength{\oddsidemargin}{-0.530in}
\addtolength{\evensidemargin}{-0.375in}
\addtolength{\textwidth}{1in}
\addtolength{\topmargin}{-0.5in}
\addtolength{\textheight}{1.1in}

\urlstyle{rm}
\raggedbottom
\raggedright
\setlength{\tabcolsep}{0in}

\titleformat{\section}{
  \vspace{-10pt}\scshape\raggedright\large
}{}{0em}{}[\color{black}\titlerule \vspace{-6pt}]

\newcommand{\resumeItem}[2]{
  \item{
    \textbf{#1}{: #2 \vspace{-2pt}}
  }
}

\newcommand{\resumeSubheading}[4]{
  \vspace{-1pt}\item
    \begin{tabular*}{0.97\textwidth}{l@{\extracolsep{\fill}}r}
      \textbf{#1} & #2 \\
      \textit{#3} & \textit{#4} \\
    \end{tabular*}\vspace{-5pt}
}

\renewcommand{\labelitemii}{$\circ$}
\newcommand{\resumeSubHeadingListStart}{\begin{itemize}[leftmargin=*]}
\newcommand{\resumeSubHeadingListEnd}{\end{itemize}\vspace{-2pt}}
\newcommand{\resumeItemListStart}{\begin{itemize}}
\newcommand{\resumeItemListEnd}{\end{itemize}\vspace{-6pt}}

\begin{document}

\begin{tabular*}{\textwidth}{l@{\extracolsep{\fill}}r}
  \textbf{\Large Junyi (Conny) Zhou} & Tel: (404) 663-5601 \\
  Boston, MA 02215 & Email: \href{mailto:junyizhou@hsph.harvard.edu}{junyizhou@hsph.harvard.edu} \\
  \href{https://github.com/JunyiZhou-Conny}{github.com/JunyiZhou-Conny} &
  \href{https://www.linkedin.com/in/junyi-zhou-270208247}{linkedin.com/in/junyi-zhou-270208247} \\
\end{tabular*}

\section{Education}
\resumeSubHeadingListStart
  \resumeSubheading
  {Harvard University}{Boston, MA}
  {M.S. in Health Data Science}{August 2025 -- December 2026}
  \resumeItemListStart
    \item Coursework: Deep Learning, Healthcare Machine Learning
  \resumeItemListEnd
  \resumeSubheading
  {Emory University, College of Arts and Sciences}{Atlanta, GA}
  {B.S. in Applied Mathematics and Statistics; Minor in Computer Informatics; GPA: 3.925}{August 2021 -- May 2025}
  \resumeItemListStart
    \item Honors: Dean's List Spring 2022, Spring 2023
    \item Coursework: Probability and Statistics, Machine Learning, NLP, Data Structures, Database Systems
  \resumeItemListEnd
\resumeSubHeadingListEnd

\section{Programming Skills}
\resumeSubHeadingListStart
  \resumeItem{Languages}{Python, C++, Java, SQL, R, JavaScript}
  \resumeItem{ML Frameworks}{PyTorch, TensorFlow, Keras, scikit-learn, XGBoost, LightGBM, Hugging Face Transformers}
  \resumeItem{Databases}{MySQL, PostgreSQL, MongoDB, Pinecone (Vector DB), Hadoop, Spark}
  \resumeItem{Tools}{Claude Code, Cursor, Git, Docker, Kubernetes, Jupyter}
  \resumeItem{Cloud}{AWS (SageMaker, Lambda, EC2, S3, Redshift), Snowflake, BigQuery}
  \resumeItem{Certifications}{AWS Certified Cloud Practitioner, AWS Certified AI Practitioner, AWS Machine Learning Associate}
\resumeSubHeadingListEnd

% Same project order and bullets as base (domain-forward)
\section{Research Projects}
\resumeSubHeadingListStart

  \resumeSubheading
    {Airway Management Simulation Chatbot}{Atlanta, GA}
    {Full-Stack Engineer, Scrum Master | Emory Pediatric Hospital}{January 2024 -- Present}
    \resumeItemListStart
      \item Architected a clinical \textbf{agentic LLM pipeline} via a \textbf{RAG} system and \textbf{OpenAI models}, minimizing hallucinations and ensuring high data privacy for simulations.
      \item Built a \textbf{HIPAA-compliant} full-stack platform (\textbf{Python/Flask}, React, MongoDB, AWS) tracking resident performance metrics via secure, real-time analytics.
      \item Implemented \textbf{Prompt Engineering} and context window management to drive dynamic scenario generation for rapid cycle deliberate practice clinical training.
      \item Led an Agile team of 6, facilitating code reviews and integrating clinician feedback to continually refine the NLP pipeline. \footnote{\url{https://d3g1qw60hz7yw7.cloudfront.net}}
    \resumeItemListEnd

  \resumeSubheading
    {Cross-Species Drug Translation using Optimal Transport \& Generative AI}{Boston, MA}
    {Graduate Researcher | Mooney Lab \& Alvarez-Melis Lab, Wyss Institute}{February 2026 -- Present}
    \resumeItemListStart
      \item Spearheaded computational frameworks using \textbf{Optimal Transport}, \textbf{VAEs}, and \textbf{diffusion models} to predict animal-to-human clinical trial translation via scRNA-seq atlases.
      \item Engineered a custom \textbf{PyTorch} CellOT implementation from scratch, modeling complex species and treatment effects directly within VAE latent spaces.
      \item Designed \textbf{scanpy} pipelines for automated cell-type labeling, integrating massive public and lab-generated scRNA-seq datasets for cancer and autoimmunity applications.
      \item Trained out-of-distribution models mapping cellular subpopulations (e.g., regulatory vs. helper T cells) across species, evaluating GENOT to optimize T cell matching.
    \resumeItemListEnd

  \resumeSubheading
    {AlphaFold Protein--Nucleic Acid Interaction Prediction}{Atlanta, GA}
    {Data Analyst | Bou-Nader Lab, Emory Dept. of Biochemistry}{February 2025 -- August 2025}
    \resumeItemListStart
      \item Engineered a pipeline integrating \textbf{AlphaFold-Multimer} and \textbf{AlphaFold3} architectures to predict complex protein--nucleic acid interactions using \textbf{MSA} data.
      \item Automated high-throughput inference on \textbf{HPC} clusters with \textbf{parallel scheduling} to efficiently process structural predictions across GPU nodes.
      \item Developed \textbf{Python} scripts (PyMOL, ChimeraX) filtering predicted structures via \textbf{pLDDT} and \textbf{PAE} scores, selecting top-performing complexes for wet-lab validation.
      \item Authored CLI tools and technical documentation to democratize advanced ML inference tools for non-technical researchers. \footnote{\url{https://github.com/JunyiZhou-Conny/AlphaFold-Bou-Nader-Lab}}
    \resumeItemListEnd

  \resumeSubheading
    {Reimplementation of Transformer Architecture}{Atlanta, GA}
    {Machine Learning Researcher | Emory Dept. of Computer Science}{October 2024 -- January 2025}
    \resumeItemListStart
      \item Implemented a Transformer in \textbf{PyTorch} from scratch, building \textbf{Multi-Head Self-Attention}, \textbf{Positional Encodings}, and \textbf{Layer Normalization} to reproduce benchmarks.
      \item Conducted extensive hyperparameter tuning and optimized the \textbf{Cross-Entropy Loss} function, achieving a 9.38 BLEU score on English-to-German translation.
      \item Benchmarked custom implementation against standard \texttt{nn.Transformer} modules to strictly validate architectural correctness and training efficiency. \footnote{\url{https://github.com/JunyiZhou-Conny/Reimplementation-of-Transformer}}
    \resumeItemListEnd

\resumeSubHeadingListEnd
\end{document}
